\documentclass[11pt]{article}
\usepackage[a4paper,margin=1in]{geometry}
\usepackage[T1]{fontenc}
\usepackage{helvet}
\renewcommand{\familydefault}{\sfdefault}
\usepackage{xcolor}
\usepackage{booktabs}
\usepackage{array}
\usepackage[font=footnotesize,labelformat=empty]{caption}
\usepackage{titlesec}
\usepackage{fancyhdr}
\usepackage{enumitem}
\usepackage{amsmath}
\usepackage{lastpage}
\usepackage[hidelinks]{hyperref}

% ---- Bayesian Capital brand palette ----
\definecolor{bcnavy}{RGB}{11,31,58}
\definecolor{bcblue}{RGB}{32,74,135}
\definecolor{bcgold}{RGB}{176,141,87}
\definecolor{bcgrey}{RGB}{90,98,112}
\definecolor{bcrule}{RGB}{210,214,220}

\titleformat{\section}{\color{bcnavy}\Large\bfseries}{\thesection}{0.6em}{}
\titleformat{\subsection}{\color{bcblue}\large\bfseries}{\thesubsection}{0.6em}{}
\titlespacing*{\section}{0pt}{1.4em}{0.5em}

\pagestyle{fancy}\fancyhf{}
\renewcommand{\headrulewidth}{0pt}
\renewcommand{\footrulewidth}{0.4pt}
\renewcommand{\footrule}{\hbox to\headwidth{\color{bcrule}\leaders\hrule height \footrulewidth\hfill}}
\fancyfoot[L]{\footnotesize\color{bcgrey}Bayesian Capital}
\fancyfoot[C]{\footnotesize\color{bcgrey}Confidential --- Internal Research}
\fancyfoot[R]{\footnotesize\color{bcgrey}Page \thepage\ of \pageref{LastPage}}

\newcommand{\kpi}[2]{\textbf{#1}\,{\footnotesize\color{bcgrey}#2}}
\setlength{\parindent}{0pt}\setlength{\parskip}{0.55em}

\begin{document}

% ---------------- Title block ----------------
\begin{titlepage}
\vspace*{1.2cm}
{\color{bcgold}\rule{\textwidth}{2pt}}\\[0.4cm]
{\color{bcnavy}\Huge\bfseries Pre-Market Signal \& Laddering Study}\\[0.35cm]
{\color{bcblue}\LARGE Findings Memorandum}\\[0.5cm]
{\color{bcgold}\rule{\textwidth}{1pt}}\\[0.6cm]
{\large\color{bcgrey} Hybrid Bayesian / Ornstein--Uhlenbeck Volatility Harvester}\\[0.2cm]
{\large\color{bcgrey} Ten-name high-beta US equity book}\\[1.2cm]
{\color{bcnavy}\large\bfseries Bayesian Capital}\\[0.15cm]
{\color{bcgrey} Quantitative Research}\\[0.15cm]
{\color{bcgrey} 29 July 2026}\\[2cm]

\fbox{\parbox{0.92\textwidth}{\color{bcnavy}\textbf{One-line conclusion.}\\ \color{black}
The model is sound as designed. The investigation yields \emph{one correctness fix}
(submit uncapped limit bids) and \emph{one optional enhancement} (Option~C laddering, for
trade volume). Pre-market VWAP and the actual open do \emph{not} improve the entry signal,
and re-optimisation does not survive out of sample.}}
\end{titlepage}

% ---------------- Executive summary ----------------
\section{Executive summary}

This memo records a systematic investigation into whether pre-market data (04:00--09:30~ET
VWAP) or the actual opening price could improve the daily Bayesian~/~OU volatility
harvester, and whether trade volume could be increased without sacrificing quality. All work
runs on a Python replica of the workbook validated to reproduce every one of the ten cached
model sheets \textbf{to the dollar}.

\begin{itemize}[leftmargin=1.4em,itemsep=2pt]
  \item \kpi{Entry signal.}{} Pre-market VWAP and the actual open, used as the price
        \emph{signal}, \textbf{do not beat the previous close}. A less-noisy signal makes the
        bid chase price and fill \emph{less} --- the smoothing is the edge, not noise to remove.
  \item \kpi{Re-optimisation.}{} Re-tuning the parameters for the open signal is an in-sample
        mirage ($+50$pp) that \textbf{does not survive out of sample}; the frozen parameters
        remain best across all ten names.
  \item \kpi{Execution fix (adopt).}{} The bid should be submitted \emph{uncapped}; a marketable
        limit fills at the open automatically. This captures the true-open behaviour the
        backtest already assumes --- a correctness fix, not an alpha change.
  \item \kpi{Trade volume (optional).}{} A Bayes-only laddered entry with a first-rung
        take-profit (\textbf{Option~C}) raises trade count \textbf{50--65\%} with
        \emph{neutral} risk-adjusted return, at a modest absolute-return cost and a modest
        softening of the OU hedge.
\end{itemize}

% ---------------- Method ----------------
\section{Scope and method}

The engine is a faithful reconstruction of each \texttt{Model} sheet: a Kalman
local-linear-trend filter (Bayes tranche) and an AR(1) mean-reverter (OU tranche), with the
50-day stop, monthly idle-cash interest and per-stock compounding. It reproduces the cached
workbook results exactly (e.g.\ NVDA: terminal \$7{,}520{,}461.83, 150.20\% annualised,
186 buys, 0 stop-loss exits), and matches annual return, buy count and stop count for all ten
names. Every candidate change is judged out of sample --- by walk-forward re-optimisation or
by sub-period consistency across names and time --- never on an in-sample fit, which always
flatters.

% ---------------- Signal study ----------------
\section{Does pre-market VWAP / the actual open improve the \emph{signal}?}

The model derives each bid from a smoothed reference: the Bayes fair value (Kalman state
through yesterday's close) and the OU forecast anchored on the previous close. We tested
replacing that reference with a ``fresher'' or ``less-noisy'' price. Execution held fixed;
frozen parameters.

\begin{table}[h]\centering\small
\renewcommand{\arraystretch}{1.15}
\begin{tabular}{lcc}
\toprule
\textbf{Signal change (all ten names)} & \textbf{avg $\Delta$ annual} & \textbf{beats baseline} \\
\midrule
OU anchor $=$ PM-VWAP        & $-8.4$pp  & 2/10 \\
OU anchor $=$ actual open    & $-10.0$pp & 2/10 \\
Bayes fair $\to$ open (gain 0.5) & $-14.6$pp & 0/10 \\
\bottomrule
\end{tabular}
\caption*{\footnotesize Fills mostly \emph{decrease}. A fresher signal makes the resting bid
track the market, so daily volatility no longer has to reach down to it.}
\end{table}

\subsection{Fair-fight re-optimisation and the all-ten tiebreaker}

Because the frozen parameters were tuned for the close signal, the open signal was granted its
own tuning (OU anchor $=$ open; Bayes fair nudged toward open with a free gain; all eleven
parameters optimised for profit plus a trade floor). In-sample this looked excellent
(NVDA 150.2\%~$\to$~200.6\%, 186~$\to$~268 buys). Out of sample it evaporated. The decisive
test re-optimised \emph{both} signals per fold across all ten names:

\begin{table}[h]\centering\small
\renewcommand{\arraystretch}{1.15}
\begin{tabular}{lc}
\toprule
\textbf{Model (avg OOS return, 30 folds)} & \textbf{return} \\
\midrule
Frozen close signal (deployed incumbent) & \textbf{52.9\%} \\
Re-optimised open signal                 & 44.9\% \\
Re-optimised close signal                & 32.8\% \\
\bottomrule
\end{tabular}
\caption*{\footnotesize The re-tuned open signal beats the incumbent in only 15/30 folds and
5/10 names --- a coin flip --- and the incumbent wins the aggregate, because the open signal
gives up its largest ground on the highest-return names. Re-optimising the close signal is
self-defeating (the recurring over-fitting lesson).}
\end{table}

\textbf{Conclusion.} As an entry signal, the previous close (via the smoothed filter) is best.
Keep it.

% ---------------- Execution fix ----------------
\section{The one correctness fix: submit uncapped bids}

The backtest caps each bid at the day's open, $\text{bid}=\min(\text{fair}-\text{buffer},\,
O_t,\,\text{peak})$. This is \emph{not} clairvoyance: it is a faithful model of limit-order
execution --- a limit buy placed above the open simply fills \emph{at} the open. It follows
that the correct live order is the \textbf{uncapped} bid
$\min(\text{fair}-\text{buffer},\,\text{peak})$, submitted before the bell; if the stock gaps
below it, the exchange fills at the open automatically.

Submitting the uncapped bid and letting the limit fill at the open reproduces the backtest
\textbf{to zero mismatches} across all 581 NVDA days. Consequently the workbook's headline
(``oracle'') returns are \emph{achievable live} --- they never required predicting the open.
The earlier notion that a pre-market proxy was needed to estimate the open cap is therefore
moot: no proxy is required. This is the single change we recommend adopting without
reservation.

% ---------------- Laddering ----------------
\section{Increasing trade volume: laddered scale-in (Option~C)}

Naive laddering (splitting each tranche's bid into rungs at increasing depth) \emph{reduces}
profit: it holds capital in reserve, starving the frequent shallow dips that drive the P\&L,
and --- when applied to both sleeves --- drives the Bayes--OU correlation from $-0.01$ to
$+0.4\ldots+0.66$, collapsing the decorrelation hedge.

Two refinements rescue it as a \emph{volume} tool:
\begin{enumerate}[leftmargin=1.6em,itemsep=2pt]
  \item \textbf{First-rung take-profit.} The whole laddered position exits at the shallowest
        rung's target, so the deeper (cheaper) rungs book their full margin rather than
        breaking even at a blended target.
  \item \textbf{Bayes-only, weighted to the shallow rung.} Laddering only the Bayes sleeve
        (OU left single-bid) preserves the hedge far better; weighting heavily to rung~1 keeps
        near-baseline deployment on normal dips.
\end{enumerate}

\textbf{Option~C} $=$ Bayes-only ladder, depths $[1.0,\,1.3,\,1.7]\times k\sigma$, weights
$[0.80,\,0.15,\,0.05]$, first-rung take-profit; OU unchanged.

\begin{table}[h]\centering\small
\renewcommand{\arraystretch}{1.15}
\begin{tabular}{lcccc}
\toprule
\textbf{Average across ten names} & \textbf{Return} & \textbf{Sharpe} & \textbf{max DD} & \textbf{Trades} \\
\midrule
Baseline (single bid) & 222\% & 4.20 & 17\% & 258 \\
Option~C ladder       & 208\% & \textbf{4.46} & 16\% & \textbf{388 ($+50\%$)} \\
\bottomrule
\end{tabular}
\caption*{\footnotesize Walk-forward (30 sub-period slices): more trades in \textbf{30/30}
slices; Sharpe neutral-to-better ($\ge$ baseline in $\approx$57\% of slices) --- confirming
this is \emph{not} the frequency-maximisation trap. Costs: $\approx$3pp per slice of absolute
return, a worst-slice lag of $\approx$20pp, and OU-hedge correlation softening to $\approx$0.19.}
\end{table}

\textbf{Assessment.} Option~C is a legitimate, out-of-sample-robust way to increase trade
volume (useful for statistical confidence in the live edge) \emph{if} a larger trade sample
and steady risk-adjusted return are valued above peak absolute return. It is a discretionary
enhancement, not a free lunch: it trades a little return and a little of the OU hedge for
$\sim$50\% more trades.

% ---------------- Recommendations ----------------
\section{Recommendations}

\begin{enumerate}[leftmargin=1.6em,itemsep=3pt]
  \item \textbf{Adopt} the uncapped-bid execution: submit the limit at
        $\min(\text{fair}-\text{buffer},\,\text{peak})$ and let it fill at the open. No open
        proxy, no PM-VWAP. This aligns live fills with the validated backtest.
  \item \textbf{Keep} the previous close as the entry signal and \textbf{freeze} the
        parameters. Neither a fresher signal nor re-optimisation survives out of sample.
  \item \textbf{Optionally adopt} Option~C laddering on the Bayes sleeve for trade volume,
        with eyes open on the two costs (a few points of return; a modestly weaker OU hedge).
        Wired into the workbook Dashboard for live use.
  \item \textbf{Do not} ladder the OU sleeve, use PM-VWAP/open as a signal, or re-optimise
        parameters --- each fights one of the two structural edges (patience, decorrelation).
\end{enumerate}

\vfill
{\color{bcrule}\rule{\textwidth}{0.4pt}}\\[-0.2cm]
{\footnotesize\color{bcgrey} Prepared by Bayesian Capital Quantitative Research. All figures
derived from a workbook-validated replica of the live model; out-of-sample results are
walk-forward or sub-period consistency tests. Past simulated performance is not indicative of
future results.}

\end{document}
